\SourceLeaf{114}%
\Article{32}
The rules for arithmetic operations on complex numbers are well known: introducing the norm reduces division to multiplication, since
\[
\frac{a+bi}{c+di}=(a+bi)\pdot\frac{c-di}{cc+dd}=\frac{ac+bd}{cc+dd}+\frac{bc-ad}{cc+dd}\pdot i
\]

The square root is extracted with the help of the formula
\[
\surd(a+bi)=\pm\left(\surd\frac{\surd(aa+bb)+a}{2}+i\surd\frac{\surd(aa+bb)-a}{2}\right)
\]
if $b$ is positive, or of the formula
\[
\surd(a+bi)=\pm\left(\surd\frac{\surd(aa+bb)+a}{2}-i\surd\frac{\surd(aa+bb)-a}{2}\right)
\]
if $b$ is negative. There is no need to dwell here on how transforming the complex quantity $a+bi$ into $r(\cos\varphi+i\sin\varphi)$ makes calculations easier.

\Needspace{17\baselineskip}
\Article{33}
We call a complex integer that can be resolved into two factors different from the units\begingroup\renewcommand{\thefootnote}{*)}\footnote{Or, equivalently, factors whose norms are greater than unity.}\endgroup\ a composite complex number; a number that admits no such factorization will be called a complex prime. It follows immediately that every real composite number is also a composite complex number. But a real prime may be a composite complex number, and this holds for the number $2$ and all positive real primes of the form $4n+1$ (except the number $1$), since they are known to be decomposable into two positive squares. For example, $2=(1+i)(1-i)$, $5=(1+2i)(1-2i)$, $13=(3+2i)(3-2i)$, $17=(1+4i)(1-4i)$, etc.

On the other hand, positive real primes of the form $4n+3$ are always complex primes. For if such a number $q$ were $=(a+bi)(\alpha+\beta i)$, we should also have $q=(a-bi)(\alpha-\beta i)$, and therefore $qq=(aa+bb)(\alpha\alpha+\beta\beta)$. But $qq$ can be resolved into positive factors greater than unity in only one way, namely into $q\times q$, so we should have to have $q=aa+bb=\alpha\alpha+\beta\beta$, \mbox{Q. E. A.}; for the sum of two squares cannot have the form $4n+3$.

\SourceLeaf{115}%
Negative real numbers evidently retain the same names as positive ones, and the same holds for pure imaginary numbers.

It therefore remains to show how to distinguish composites from primes among mixed imaginary numbers. This is done by the following

\textsc{Theorem.}\quad\emph{Every mixed imaginary integer $a+bi$ is a complex prime or a composite number according as its norm is a real prime or a composite number.}

\emph{Proof.} I. Since the norm of a composite complex number is always composite, a complex number whose norm is a real prime must evidently be a complex prime. \mbox{Q. E. P.}

II. If, however, the norm $aa+bb$ is composite, let $p$ be a positive real prime dividing it. We must now distinguish two cases.

1) If $p$ has the form $4n+3$, we know that $aa+bb$ cannot be divisible by $p$ unless $p$ divides both $a,b$, so that $a+bi$ will be composite.

\begingroup\postdisplaypenalty=10000\relax
2) If $p$ does not have the form $4n+3$, it can certainly be decomposed into two squares. We therefore put $p=\alpha\alpha+\beta\beta$. Since
\[
(a\alpha+b\beta)(a\alpha-b\beta)=aa(\alpha\alpha+\beta\beta)-\beta\beta(aa+bb)
\]
and this is therefore divisible by $p$, $p$ must divide one of the factors $a\alpha+b\beta$, $a\alpha-b\beta$. Moreover, since
\[
(a\alpha+b\beta)^2+(b\alpha-a\beta)^2=(a\alpha-b\beta)^2+(b\alpha+a\beta)^2=(aa+bb)(\alpha\alpha+\beta\beta)
\]
and this is therefore divisible by $pp$, in the first case $b\alpha-a\beta$, and in the second $b\alpha+a\beta$, must also be divisible by $p$. Thus in the first case
\[
\frac{a+bi}{\alpha+\beta i}=\frac{a\alpha+b\beta}{p}+\frac{b\alpha-a\beta}{p}\pdot i
\]
will be a complex integer, while in the second
\[
\frac{a+bi}{\alpha-\beta i}=\frac{a\alpha-b\beta}{p}+\frac{b\alpha+a\beta}{p}\pdot i
\]
will be an integer. Since the given number is therefore divisible by $\alpha+\beta i$ or by $\alpha-\beta i$, and the norm of the quotient $=\frac{aa+bb}{p}$ differs from unity by hypothesis, $a+bi$ is evidently a composite complex number in either case. \mbox{Q. E. S.}\endgroup

\SourceLeaf{116}%
\Article{34}
The following four kinds therefore exhaust the entire range of complex primes:

\begin{list}{}{\setlength{\leftmargin}{1.7em}\setlength{\labelwidth}{1.5em}\setlength{\labelsep}{.2em}}
\item[1)] the four units, $1,+i,-1,-i$, which, however, we shall usually leave understood to be excluded when discussing prime numbers;
\item[2)] the number $1+i$ with its three associates $-1+i$, $-1-i$, $1-i$;
\item[3)] the positive real primes of the form $4n+3$ with their three associates each;
\item[4)] the complex numbers whose norms are real primes of the form $4n+1$ greater than unity. To each given norm of this kind there always correspond eight complex primes and no more, since such a norm can be decomposed into two squares in only one way.
\end{list}

\Article{35}
Just as real integers are divided into even and odd, and the former in turn into those divisible by four and those that are twice an odd number, an equally essential distinction presents itself among complex numbers. They are either

not divisible by $1+i$, namely numbers $a+bi$ in which one of the numbers $a,b$ is odd and the other even;

or divisible by $1+i$ but not by $2$, whenever both $a,b$ are odd;

or divisible by $2$, whenever both $a,b$ are even.

The numbers of the first class may suitably be called \emph{odd complex numbers}, those of the second \emph{half-even}, and those of the third \emph{even}.

A product of several complex factors will always be odd when all the factors are odd; half-even when one factor is half-even and the others odd; and even when the factors include at least two half-even ones or at least one even one.

The norm of any odd complex number has the form $4n+1$; the norm of a half-even number has the form $8n+2$; finally, the norm of an even number is the product of a number of the form $4n+1$ and the number $4$ or a higher power of two.

\Needspace{18\baselineskip}
\Article{36}
Since the relation between four associated complex numbers is analogous to that between two opposite real numbers (that is, numbers equal in absolute value and bearing opposite signs), and the positive one of these is usually and rightly regarded as primary,\ \SourceLeaf{117}%
the question arises whether a similar distinction can be established among four associated complex numbers, and whether it should be considered useful. To decide this, we must consider that the principle of distinction must be such that the product of two numbers regarded as primary among their respective associates always becomes primary among its own associates. But we soon become convinced that no such principle exists at all unless the distinction is restricted to integers; indeed, a \emph{useful} distinction must be limited to odd numbers. For these, however, the proposed aim can be reached in two ways. Namely,

I. The product of two numbers $a+bi,a'+b'i$ such that $a,a'$ have the form $4n+1$ and $b,b'$ are even will have the same property: the real part will be $\equiv1\ (\mathrm{mod.}\ 4)$, and the imaginary part even. And it is easy to see that among four associated odd numbers only one has that form.

II. If a number $a+bi$ is such that $a-1$ and $b$ are either both divisible by four or both twice an odd number, its product with a complex number of the same form will have the same form; and it is easy to see that only one of four associated odd numbers has this form.

\Needspace{11\baselineskip}
Of these two principles, which are almost equally suitable, we shall adopt the latter: among four associated odd complex numbers we shall regard as primary the one congruent to the positive unit modulo $2+2i$. This will let us state several notable theorems more concisely. For example, the complex primes $-1+2i$, $-1-2i$, $+3+2i$, $+3-2i$, $+1+4i$, $+1-4i$, etc., are primary, as are the real numbers $-3$, $-7$, $-11$, $-19$, etc., which must evidently always be taken with a negative sign. The conjugate of an odd primary complex number will also be primary.

For half-even and even numbers in general, a similar distinction would be too arbitrary and of little use. Of the associated primes $1+i$, $1-i$, $-1+i$, $-1-i$, we can indeed choose one as primary in preference to the others, but we shall not extend such a distinction to composites.

\Article{37}
If some factors of a composite complex number are themselves composite, and these in turn are resolved into their factors, we evidently arrive eventually at prime factors; that is, every composite number can be resolved into prime factors. If any of these are not primary, replace each by\ \SourceLeaf{118}%
the product of its primary associate and $i,-1$ or $-i$. In this way, every composite complex number $M$ can evidently be reduced to the form
\[
M=i^\mu A^\alpha B^\beta C^\gamma\dotsfour
\]
where $A,B,C$, etc., are distinct primary complex primes, and $\mu=0,1,2$ or $3$. Concerning this factorization, we have the theorem that it can be carried out in only one way. On a passing consideration this theorem might seem evident in itself, but it certainly needs a proof. The following paves the way to that proof.

\textsc{Theorem.}\quad\emph{The product $M=A^\alpha B^\beta C^\gamma\dotsfour$, where $A,B,C$, etc., denote distinct primary complex primes, cannot be divisible by any primary complex prime that does not occur among $A,B,C$, etc.}

\emph{Proof.} Let $P$ be a primary complex prime not contained among $A,B,C$, etc., and let $p,a,b,c$, etc., be the norms of the numbers $P,A,B,C$, etc. It follows readily that the norm of $M$ will be $=a^\alpha b^\beta c^\gamma$, etc.; hence, if $M$ were divisible by $P$, this number would have to be divisible by $p$. Since the individual norms are either real primes (from the sequence $2,5,13,17$, etc.) or squares of real primes (from the sequence $9,49,121$, etc.), this evidently cannot happen unless $p$ is identical with one of the norms $a,b,c$, etc. We therefore suppose $p=a$. But since $P,A$ are distinct primary complex primes by hypothesis, it is easy to see that these conditions cannot hold together unless $P,A$ are conjugate imaginary complex numbers, and consequently $p=a$ is an odd real prime (not the square of a prime). We therefore suppose $A=k+li$, $P=k-li$. Hence, extending the notion and notation of congruence to complex integers, we have $A\equiv2k\ (\mathrm{mod.}\ P)$, from which it follows readily that
\[
M\equiv2^\alpha k^\alpha B^\beta C^\gamma\dotsfour\ (\mathrm{mod.}\ P)
\]
Therefore, if $M$ is assumed divisible by $P$, then
\[
2^\alpha k^\alpha B^\beta C^\gamma\dotsfour
\]
will also be divisible by $P$, and consequently the norm of this number, which is
\[
=2^{2\alpha}k^{2\alpha}b^\beta c^\gamma\dotsfour
\]
will be divisible by $p$. But since $2$ and $k$ are certainly not divisible by $p$, it follows\ \SourceLeaf{119}%
that $p$ must be identical with one of the numbers $b,c$, etc.; \mbox{let, for example, $p=b$.} From this, however, we conclude that either $B=k+li$ or $B=k-li$, that is, either $B=A$ or $B=P$, both contrary to the hypothesis.

The other theorem, that factorization into primes can be carried out in only one way, follows very easily from this theorem, by reasoning entirely analogous to that which we used for real numbers in the \emph{Disquisitiones Arithmeticae} (Article 16). It would therefore be superfluous to dwell on it here.

